\section{Our Release Protocol}
\label{sec:protocol-evaluation}
Our protocol assumes a central public agency that maintains one
authoritative dependency graph for the entire dataset under its
responsibility. The graph covers the original data, derived dataset
versions, training computations and model releases. This setting
corresponds to an agency allowing public access to approved versions
of its data for model training. The agency audits proposed releases
and maintains a common privacy account across users and projects.

At ingestion, direct identifiers are removed and a separate restricted
linkage boundary supplies stable pseudonyms. Pseudonymization does not
make records anonymous or prevent inference from other fields. External
users receive only accounted private releases; both data and models enter
the pooled history. The ingestion and vault requirements remain specified
architecture, not deployment evidence (Appendix~\ref{app:stage-zero}).

\begin{figure*}[!t]
\centering
\includegraphics[width=\textwidth]{figures/central-audit-graph.pdf}
\caption{Agency dataset and release history. Both ingestion paths
remove direct identifiers before analysis; a separate vault controls
identity lookup, without preventing inference from other fields.
Registering $B$ creates a
new main-dataset version; every derived dataset links to its source
version. Models $R_1,R_2$ reuse $A_1$, while $R_4$ combines $A_2,A_3$.
Solid arrows show dependencies; the dashed line shows shared people.
Prior disclosures and charges persist. New data require reviewed
membership, contribution bounds and privacy evidence. Illustrated
external training uses only accounted DP releases and fixed context.
The graph can register changing datasets, but the finite continuation
analyzed here uses one unchanged data state. The vault and ingestion
controls are specified, not deployment evidence.}
\label{fig:release-protocol}
\end{figure*}

\subsection{Register data, computations and every release}
\label{sec:central-graph}
Write the agency's directed dependency graph as
$G_t=(\mathcal D_t\cup\mathcal C_t\cup\mathcal R_t,\mathcal E_t)$.
Dataset nodes $\mathcal D_t$ identify immutable versions and their
protected units. Each derived dataset links to the relevant version of
the agency's main dataset. An addition creates a registered successor
version and updated membership records, while retaining the earlier
version and its release history (Figure~\ref{fig:release-protocol}).
Computation nodes $\mathcal C_t$ identify protected
accesses, transformations, randomness dependencies and privacy evidence.
Release nodes $\mathcal R_t$ identify each shared private dataset,
exported model package or other observable disclosure, including a
refusal without a model, and
its authorization. Edges state which inputs produced each result;
every retraining or deletion update adds a new version. Re-delivery of
a saved package refers to its existing computation, rather than silently
creating another draw under the same identity.

The agency also maintains confidential unit-membership links across
dataset nodes. Separate branches need not contain different people.
Dataset names and recipients do not alter this account. Unresolved
overlap requires a conservative bound or withheld authorization. Pairwise
overlap counts do not establish each person's complete exposure; unit
membership or a sound higher-order bound is required. These confidential
links are fixed in our adjacency and are not public exports.

\subsection{Central audit and authorization}
\label{sec:central-audit}
For a proposed release, the agency traces its ancestors and the prior
release history. It checks three matters before authorization.

\paragraph{Complete scope and credible execution.}
Reconcile declared inputs, protected units and existing releases with
the controlled data and export inventory. Include protected labels,
preprocessing, selection, validation and visible audit outcomes.
Bind the proposed package to the audited data version and computation.
A submitted graph or artifact hash alone does not establish truthful
execution. Appendix~\ref{app:central-audit} separates these obligations
from the local prototype's evidence.

\paragraph{A valid joint privacy argument.}
Audit the whole dependency path, not a model's attack score. Reuse is
free only when outputs are post-processing of an already-accounted
private source. Fresh protected access needs conditional composition
or an explicit joint-channel proof; correlated marginal guarantees
cannot simply be added. These arguments must hold over the declared
neighboring datasets and permitted histories, not only the realized
model bytes. Charge each certified source for each affected
unit once, not once per descendant model. The finite constructions
below provide particular joint arguments. Missing privacy evidence
blocks release; an audit cannot retrospectively repair an earlier
invalid disclosure.
The charge must cover the unit's full contribution to each source.

\paragraph{Task utility.}
Freeze the package, prediction inputs and task requirements before
independent evaluation. Passing a task requirement and passing privacy
checks are separate requirements; a useful model or unsuccessful attack
does not certify DP. Changed data, including deletion, require a new
mechanism argument; registration alone does not make a fixed-data
extension applicable.

\subsection{Extend a known finite release channel}
\label{sec:commitment}
Suppose the pooled earlier output $O$ has known conditional law
$P(o\mid s)$. It includes disclosures to every covered recipient;
the finite construction can treat this bundle as one old channel.
A new workload requires $r$ independent uniform binary coordinates
$Z_1,\ldots,Z_r$, independent of the old protected input $s$.
The agency may access the data again, but must retain the actual old
output. The complete pair of old output and new decisions must satisfy
pure $E$-DP under replacement of $(s,Z)$. We optimize minimum new-task
accuracy, averaging over a declared public prior on $s$.

This construction treats $(s,Z)$ as one protected unit and compares
every pair of its values. Exposing its separately randomized row outputs
gives the local-DP observation model, even when the agency samples them.
It answers how to extend an irrevocable finite disclosure, rather than
which mechanism the agency should choose before any disclosure exists.
For population learning, a central computation that never exposes
row-level responses is a necessary comparator. The model-history
formulation below handles a different, general record-neighbor graph;
its extension principle does not require a local old release.

Put $N=2^r$ and $\lambda=e^E$. Symmetry groups new error vectors by
their number $h$ of incorrect coordinates. For each old output choose
a probability floor
\[
 \frac{\max_sP(o\mid s)}{N\lambda}
 \le m_o\le \frac{\min_sP(o\mid s)}N.
\]
For each old state, assign that floor to every error vector and fill
remaining probability in increasing order of $h$, up to $\lambda m_o$
per vector. The optimum occurs at an endpoint or one of finitely many
scalar breakpoints. Proposition~\ref{prop:finite-commitment} gives the
formula, boundary cases and proof. For binary randomized response (RR)
as the old channel, the calculation takes $O(r^2)$ arithmetic operations.
It requires the complete old channel, not merely its privacy scalar.
It is not a tractable characterization of arbitrary trained weights.

The one-new-coordinate formula recovers prior work
\citep{yamamoto2024privacyoptimized}. For several coordinates the
proper comparison is an optimized joint mechanism for the new vector,
independent of the old response, not only coordinate-wise RR.
\subsection{Use the actual retained model history}
\label{sec:weights-main}
Let $P_D(m)$ be the joint law of the pooled old history, including
weights and other data-dependent disclosures to any covered recipient.
Separate models' marginal laws do not determine this joint law.
An extension $Q_D(m,a)$
must satisfy
\[
 \begin{aligned}
 \sum_a Q_D(m,a)&=P_D(m),\\
 Q_D(m,a)&\le e^E Q_{D'}(m,a)\quad(D\sim D').
 \end{aligned}
\]
Given the actual $m$, the agency can realize a feasible finite extension
by sampling $a$ from $Q_D(m,a)/P_D(m)$ on its positive support. This
requires protected data access and knowledge of the conditional laws;
it is not free post-processing of $m$.

Summing the joint inequality over $a$ shows that the old law must itself
satisfy the requested bound. A continuation cannot repair a nonprivate
earlier release. For example, a deterministic old output that differs on
neighbors already gives an event of probability one versus zero and
cannot admit a finite pure-DP extension. Knowing a law is therefore
necessary but not sufficient for feasibility.

Protecting a full private training representation is sufficient but can
be unnecessarily restrictive when only particular weights were exported
and no recipient obtained that representation.
If private training data releases are shared, they remain part of the old
disclosure history; a model-only comparison cannot omit them.
Appendix~\ref{sec:actual-model-history} proves a strict separation in a
three-record learning problem and a recoverability condition for equality.
An omitted disclosure invalidates the comparison. The result is not a
budget refund or a general model analyzer; ordinary composition remains
a conservative baseline.

For a fixed $q$-value record alphabet, suppose the old channel and the
utility objective are invariant to record permutations. Averaging a
feasible extension over those permutations preserves its old marginal,
record-level privacy and utility. The channel can then depend only on
the $q$ category counts, with $\binom{n+q-1}{q-1}$ states; a neighboring
state moves one record between categories. Proposition~\ref{prop:count-reduction}
proves the reduction and gives the necessary multiplicity weights.
For our two-bit records and two binary model outputs, this gives
$4\binom{n+3}{3}$ variables instead of $4^{n+1}$.
It requires invariance of the complete earlier disclosure, not merely
an informal claim that a model ignores row order. Count aggregation
removes row-order redundancy, but still covers every possible count
state, rather than just the observed count vector. Both the record count
$n$ and alphabet size $q$ limit tractability.

With uncertain public utility laws, optimize a mixture or worst-case
regret over a declared set; the privacy constraints stay unchanged.
Choosing a task from protected data instead needs a separate charge.
Appendix~\ref{app:uncertain-workloads} proves the bound for a fixed menu
of joint channels and a private selector after the old export.



\subsection{Check whether private adaptation can help}
\label{sec:information-value}
Before reserving a continuation, compare its potential utility with
public prediction and reuse. For the declared average linear objective
$U(Q)=c_0+\sum_{D,m,a}c_{Dma}Q_D(m,a)$, define
\begin{align}
 U_{\rm reuse}&=c_0+\sum_m\max_a\sum_D P_D(m)c_{Dma},\nonumber\\
 U_{\rm full}&=c_0+\sum_{D,m}P_D(m)\max_a c_{Dma}.
 \label{eq:information-value}
\end{align}
The first is the best prediction rule using only the old output.
The second allows the rule to see the complete private state without
any privacy restriction. If $U_{\rm full}-U_{\rm reuse}<\Delta$ for a
prespecified required gain $\Delta>0$, no private continuation in this
family can meet that gain. This check needs no LP, although it still
enumerates the finite public problem. Otherwise a certified private
upper bound can sharpen the check. Passing either check only leaves
improvement possible; it does not authorize release.

When the check rules out adaptation, the agency can construct the
old-only rule directly and certify its joint history at the old law's
bound. Choosing it before reservation can avoid an unnecessarily large
charge. This is tighter prospective accounting, not stronger privacy
for an identical channel or a refund after release. The conclusion is
conditional on the public utility law and model class; it says nothing
about arbitrary unknown populations (Appendix~\ref{app:information-value}).

\subsection{Commit the release and audit subsequent changes}
\label{sec:workload-extension}
\label{sec:implemented-assurance}
Reserve the required allowance before the approved protected computation
and commit the audited release record before external access. Concurrent
proposals must fit the common allowance jointly. Only the authorized
package may leave the controlled boundary. Changes to data, code,
private inputs, noise relationships or disclosed metadata require a new
audit. Reconcile actual exports with the graph; preserve prior release
nodes after revocation or deletion. These are central control requirements,
not evidence of deployed multi-agency enforcement.

Our bounded instantiation admits independent private sources, exact reuse
and a checked finite-model continuation on fixed registered memberships.
Before accepting protected inputs, it freezes the complete public problem,
model meanings, selection kernel and the actual certified channel menu.
The optimizer uses only that public specification. Exact checks bind each
candidate's old marginal, adjacency, requested odds and objective. Merely
fixing the optimization problem is insufficient: privately choosing between
two valid solutions could reveal data. The whole history allowance is
reserved before the old draw. After any accounted
task selection, the executor samples the new model conditional on the
saved old output and unchanged confidential counts. One extension slot
and immutable committed models prevent retries or other recipients from
obtaining uncharged draws. At total odds $3$, the conservative rational
charge is $11/10>\log3$; unused or revoked reservations are not refunded.
This covers a prospectively frozen menu on the same data state, not arbitrary later
tasks or changed records. Appendix~\ref{app:finite-enforcement} gives
the execution argument.

For larger rosters, a fixed public partition can assign each privacy unit
to exactly one block, with independent mechanism randomness across
blocks. Each block mechanism uses only that block's protected values
and fixed public context. The executor applies the same certified finite
template to each block and exports their heads as an ensemble. A changed unit affects
only one complete block history; the pooled bound is the largest block
bound, not their sum. Voting is post-processing. All contributions from
one protected unit must remain in that block, and new partitions or
overlapping histories cannot reset its account. This uses standard
parallel composition, not optimal joint learning on the full roster
(Appendix~\ref{app:block-execution}).

Repeated verification of identical public certificates and plan templates
can be cached. Cache identity includes the full canonical inputs,
verification limits, model meanings and current mathematical charge.
Source binding, reservation, revocation and release authority are checked
against current state on every request. A cache hit is never itself an
authorization.

Other correlated-noise declarations remain unsupported.
Appendix~\ref{app:central-gate-proof} proves the account invariant under
truthful mechanisms, complete mediation and durable transitions. Graph
and ledger representations give equivalent admission under the same evidence.

Evaluate frozen candidates with independent public evidence or an
accounted private procedure, including selection. Identify what evidence
could resolve a threshold-crossing uncertainty; failure of tested
candidates is not general infeasibility.

